\cotefile{83}
\scanpage{178}
\unsure{les homomorphismes voulus|l'homomorphisme voulu} de
\[
\deleted{\unsure{\mathop{*}_{i}|\bigotimes_i|*}}\, E_1\mathbin{\underset{\lambda}{*}}E_2\quad\text{dans}\quad E_1^\circ * E_2^\circ .
\]

Par passage à la limite etc., on en conclut l'injectivité si les $F_i$ sont plats.

Il faut définir un $\lambda$-scindage \inserted{\uline{canonique}} sur $* E_i$, donc il faut définir
\[
\mathop{*}^{i}E_i \longrightarrow \bigotimes E_i^\circ
\]
qui sur $\bigotimes E_i^\circ$ prenne la valeur $\lambda\,\mathrm{id}$. Il faut donc trouver un $\lambda$-accouplement can. à valeurs dans $\bigotimes E_i^\circ$.

On songe \deleted{tout} de suite à
\[
\begin{aligned}
f=\bigotimes u_i &: \bigotimes E_i\longrightarrow\bigotimes E_i^\circ\\
f(x_1,\ldots x_n)&=u_1(x_1)\otimes\ldots\otimes u_n(x_n)
\end{aligned}
\]
% slanted
\margin{mais il faudrait remplacer $\bigotimes E_i$ par son image -- ou bien ne pas \unsure{exiger|croire} dans les données $F_i\xrightarrow{\unsure{R_i|\beta_i}}E_i$ que celles-ci \unsure{sont|soient} injectives mais \unsure{seulement} donner un $u_i:E_i\mathbin{\unsure{-|\to}}F_i$ avec $u_i R_i=\lambda\,\mathrm{id}_{F_i}$.}
\[
\begin{aligned}
 f(x_1+\alpha_1,x_2\ldots x_n)&=(u_1(x_1)+\lambda\alpha_1)\otimes u_2(x_2)\ldots\otimes\unsure{u_i|u_n}(x_n)\\
 &=f(x_1,\ldots x_n)+\underset{\deleted{\unsure{g(x_2)}}}{\underline{\lambda\alpha_1}}\otimes u_2(x_2)\ldots u_n(x_n)
\end{aligned}
\]
prendre donc comme application associée
\[
 f^\circ(\alpha_1,\ldots,\alpha_n)=\lambda\alpha_1\otimes\ldots\otimes\alpha_n
\]
de sorte que les termes \unsure{en $\alpha$'s|non nuls} précédents \unsure{s'écrivent|s'écrive}
% archive sheet number: 177
\scanpage{179}
Ceci dit, pour définir \inserted{un} $f$, \deleted{$\unsure{f^\circ|f_0}$} correspondant à un $f^\circ$ déjà donné, on peut se le donner arbitrairement \inserted{(soit $g$)} sur $\prod F_i$ (comme application multilin.)

\deleted{On aura alors}

Faisons le calcul pour deux facteurs -- On posera (notant $x_i\in F_i$, $\alpha_i\in E_i^\circ$)
\[
\begin{aligned}
 f(\underbrace{x_1+\alpha_1}_{X_1},\underbrace{x_2+\alpha_2}_{X_2})
 &=g(x_1,x_2)+\deleted{\unsure{f^\circ(x_1,b_2(x_2))}}f^\circ(\alpha_1,b_2(x_2))\\
 &\quad+f^\circ(b_1(x_1),\alpha_2)+
 \underbrace{\lambda f^\circ(\alpha_1,\alpha_2)}_{\substack{\Vert\\f^\circ(u_1(\alpha_1),\alpha_2)\\=f^\circ(\alpha_1,u_2(\alpha_2))}}
\end{aligned}
\]

Il suffit de vérifier que ça satisfait aux deux conditions. \unsure{Pr ex|P. ex.} si on remplace $\alpha_1$ \deleted{par} $\alpha_1+\alpha'_1$, on obtient
\[
\begin{aligned}
f(X_{1\deleted{\unsure{2|,}}}+\alpha'_1,X_2)&=f(X_1,X_2)+f^\circ(\alpha'_1,\underline{b_2(x_2)}).\\
&\quad+\lambda f^\circ(\alpha'_1,\alpha_2)
\end{aligned}
\]
et on veut
\[
 = f(\unsure{x_1|X_1},X_2)+f^\circ(\alpha'_1,\underbrace{u_2(X_2)}_{b_2(x_2)+\lambda\alpha_2})\qquad \text{OK.}
\]

Prenons pour $B$ le produit tensoriel $E_1^\circ\otimes E_{\unsure{2|3}}^\circ$, pour $f^\circ$ l'accouplement can., et $g=0$, on trouve
% archive sheet number: 178
\scanpage{180}
\[
 f^\circ(\alpha_1,u_2(x_2),\ldots,u_n(x_n)).
\]
% slanted
\margin{\uline{NB} Si les \deleted{\unsure{Ext.}} $E_i$ sont munis de scindages, et leur structure de $\lambda$-triv. définie par des $b_i:F_i\mathbin{\unsure{-|\to}} E_i^\circ$, alors la $\lambda$-trivialisation de $*E_i$, scindé de \unsure{façon} can., est donnée par $\bigotimes b_i:\bigotimes F_i\mathbin{\unsure{-|\to}}\bigotimes E_i$.}

Ainsi, on a une opération interne $*$ dans la catégorie des extensions $\lambda$-scindées \inserted{(à quotients $F_i$ plats, au besoin)}, et en particulier dans celle des torseurs $\lambda$-scindés (car si les $F_i$ sont $\simeq k$, $\bigotimes F_i\simeq k$ canon.$^{\mathrm t}$).

\deleted{\unsure{Q}} Outre \unsure{sa|la} \unsure{donnée|condition} de commutativité, d'associativité (à vérifier leurs compatibilités !), a-t-on une unité ? Il faudrait pour le moins, \deleted{que} pour une unité, \unsure{q}
\diagram{\begin{tikzcd}[ampersand replacement=\&,column sep=small]
0\arrow[r,-]\& E^\circ\arrow[r,-]\& E\arrow[r,-]\& F\arrow[r]\& 0
\end{tikzcd}}
que $E^\circ$ \deleted{soit} $F$ soient unités pour le $\otimes$ ordinaire, donc \deleted{inversibles} \inserted{$\simeq k$}.
\diagram{\begin{tikzcd}[ampersand replacement=\&,column sep=small]
0\arrow[r]\&k\arrow[r]\&E\arrow[r]\&k\arrow[r]\&0
\end{tikzcd}}

\deleted{\unsure{Supp.|Une} tels} Choisissons scindage, \deleted{alors} la $\lambda$-triv. est définie par
\[
 \deleted{\unsure{k\to k}}\unsure{b_0|\beta_0|t}:k\longrightarrow k\qquad\text{i.e. }\beta_0\in k
\]

J'ai envie de prendre $\beta_0=1$. Prenons plus généralement $\beta$ quelconque, je dis qu'alors $E_\beta\mathbin{\underset{\lambda}{*}}E_2^{\unsure{+}}\simeq{}^{\beta}E_2$, muni de la $\lambda$ trivialisation $\beta u_2$ déduite de celle de $E_2$
% archive sheet number: 179
\scanpage{181}
Ce foncteur $E\mathbin{\unsure{-|\mapsto}}{}^\beta\deleted{\unsure{V}}E$ \deleted{\unsure{se|$\to$}} est une équivalence de cat ssi $\beta$ est inversible dans $k$ -- Et si $\beta=1$, on retrouve $\mathord{}^1E=E$.

En termes intrinsèques, pour \deleted{que} un $M$ $\lambda$-trivialisé, les conditions suivantes sont équivalentes

1) $M\mathbin{\underset{\lambda}{*}}\,?$ est une équivalence sur la \unsure{Cat.|cat.} des $\lambda$-trivialisables

2) $M^\circ\mathbin{-}M$ injectif, $\underset{L''}{M^\circ}$ et $\underset{L'}{M/M^\circ}$ inversibles, $M\xrightarrow{u}M^\circ$ est surjectif
(en termes d'\deleted{$M/M^\circ$}un scindage cela signifie que l'idéal somme des \deleted{\unsure{images|id. pr.}} de $L'\otimes L^{-1}\xrightarrow{\beta} k$ et $\lambda k$ est $k$)

3) $M$ inv. pour $\mathbin{\underset{\lambda}{\deleted{\unsure{\otimes}}*}}$
% archive sheet number: 180
\scanpage{183}
\mspage{(I bis)}
\hfill \unsure{Mai} 1986\par
\hfill Juillet 1986

$\unsure{A|\mathcal A}$ alg. d'Az. de rang 4 sur $k$, on veut expliciter sa structure multiplicative en termes de la forme $\det$, de l'élément $1$, et de la structure \uline{spéciale} orthogonale $e\in\bigwedge^4 A$ \deleted{associée} \inserted{can.} associée à $\det$ ($e\otimes e=\mathrm{dis}(\det)$). Comme on connait la forme $\operatorname{Tr}xy$ par la formule
% slanted
\margin{\uline{NB} $\operatorname{Tr}$ \unsure{connue|connu} en termes de $\det$ et de $1$, comme $\operatorname{Tr}(x)=\psi_{\det}(x,1)$}
\[
 \det(x\unsure{+|-}y)=\det x+\det y+\operatorname{Tr}x\operatorname{Tr}y-\operatorname{Tr}xy
\]
i.e.
\[
 \psi_{\det\deleted{\unsure{(x,y)}}}(x,y)=\operatorname{Tr}x\operatorname{Tr}y-\underbrace{\operatorname{Tr}xy}_{\varphi(x,y)},
\]
\deleted{pour} et que celle-ci est non dég., pour connaître $xy$, il suffit de connaître
\[
 \operatorname{Tr}xyz\qquad\forall z\in A.
\]

Donc je cherche une formule pour $\operatorname{Tr}xyz$. Comme $\operatorname{Tr}(xy)z$ et $\operatorname{Tr}(yx)z$ ne sont pas égaux en général, i.e la structure mult. \inserted{de $A$} n'est \uline{pas} commutative
\deleted{\unsure{(Ainsi|(alors}} On sent que $\det,1$ sans $e$ ne suffisent à expliciter $\operatorname{Tr}xyz$ car ils ne déterminent la structure multiplicative que mod passage à l'inverse, i.e. mod l'involution
\[
 x\longmapsto x'=\operatorname{Tr}(x).1-x.
\]

On note aussi que $\operatorname{Tr}x\unsure{,}yz$ ne dépend que du triple $\{x,y,z\}$ et d'un ordre circulaire sur \unsure{l'ens|l'un|l'ens.} des trois termes. On a, \deleted{pour} si
\[
\begin{aligned}
\tau_{\unsure{,}}\deleted{\unsure{=}}&=\operatorname{Tr}xyz\quad(=\operatorname{Tr}yzx=\operatorname{Tr}zxy)\\
\tau'&=\operatorname{Tr}(zyx)(=\operatorname{Tr}(yxz)=\operatorname{Tr}(xzy))
\end{aligned}
\]
% archive sheet number: 182
\scanpage{184}
\margin{\diagram{\resizebox{\linewidth}{!}{$\begin{array}{lll}
\operatorname{Tr}x=b&\operatorname{Tr}y=b'&\operatorname{Tr}z=\unsure{b''|b'}\\
\det x=c&\det y=c'&\det z=\unsure{c''}\\
\operatorname{Tr}yz=t&\operatorname{Tr}zx=t'&\operatorname{Tr}xy=t''
\end{array}$}}}

la formule
\[
\begin{aligned}
 \tau+\tau'&=\operatorname{Tr}(xy+yx)z\qquad\text{où}\\
 xy+yx&=\underset{b'}{\underline{\operatorname{Tr}(y)}}x+\underset b{\underline{\operatorname{Tr}(x)}}y+
 (\underset{t''}{\underline{\operatorname{Tr}(xy)}}\deleted{\unsure{.1}}-\underset b{\underline{\operatorname{Tr}x}}\ \underset{b'}{\underline{\operatorname{Tr}y}})1
\end{aligned}
\]
donc
\[
\begin{aligned}
 &\deleted{(xy)z+}\\
 &(xy+yx)z=b'xz+byz+(t''-bb')z\qquad\text{d'où}\\
 &\boxed{\tau+\tau'=b^{\deleted{\unsure{\prime}}}t^{\deleted{\unsure{\prime}}}+b't'+b''t''-bb'b''}
 \overset{\mathrm{déf}}{=}\Sigma(b,b',b'',t,t',t'')
\end{aligned}
\]

On devine que $\tau\tau'$ \unsure{sera|ser.} également un polynôme en $b,b',b'',c,c',c'',t,t',t''$, de sorte que $\tau,\tau'$, \unsure{s'expliciteront|s'expliciteraient|s'explicitent} (à permutation près) comme racines d'un polynôme \inserted{unitaire} du second degré, dont les \unsure{coeff.|coeffts} $B$ ($=\tau+\tau'$) et $C$ ($=\tau\tau'$) sont des polynômes en les variables $b^{(i)},c^{(i)},t^{(i)}$.

\deleted{Mais plutôt que de chercher une expression pour $\tau\tau'$, je vais essayer de les \inserted{($\tau$ et $\tau'$)} déterminer directement en termes de $\det,1$ \uline{et de $e$}. Je suis guidé par le fait que $xy$ et $yx$ sont opposés mod $1,x,y$ -- donc \uline{si $z$ est ortho} à $1,x,y$, \deleted{det} \inserted{(pour $\unsure{\varphi|\rho}$)} on aura
\[
 \operatorname{Tr}(xyz)=-\operatorname{Tr}(yxz).
\]
Pour se ramener à ce cas, on construit $\alpha,\beta,\gamma\in k$ tels que
\[
 z'=z-\alpha1-\beta x-\gamma y
\]
soit orth. à $1,x,y$, i.e.
\[
\begin{aligned}
 b''-2\alpha-b\beta-b'\gamma&=0\\
 t'-b^{\deleted{\unsure{\prime}}}\alpha+(2c-b^2)\beta-t''\gamma&=0\\
 t^{}-\deleted{\unsure{(b^2-2c)\alpha}}b'\alpha\unsure{-|+}t''\beta+(2c'-b^{\prime 2})\gamma&=0
\end{aligned}
\]
i.e
\[
 \left\{\begin{aligned}
 2\alpha+b\beta+b'\gamma&=b''\\
 b\alpha+(b^2-2c)\beta+t''\gamma&=t'\\
 b'\alpha+t''\beta+(b^{\prime 2}-2c')\gamma&=t
 \end{aligned}\right.
 \quad\text{matrice}\quad
 \begin{pmatrix}2&b&b'\\b&b^2-2c&\unsure{t''|t'}\\b'&t''&b^{\unsure{\prime\prime|\prime 2}}-2c'\end{pmatrix}
\]}
% archive sheet number: 183
\scanpage{185}
avec \deleted{\unsure{une}}
\[
 (1\wedge x\wedge y\wedge z)\otimes e^{-1}=2\varepsilon_0\qquad
 \text{où}\quad\varepsilon_{\unsure{0|-}}=\unsure{[|(}\alpha\wedge\alpha'\wedge x_0\wedge y_0)\underset{\in\pm1}{\unsure{\otimes}e^{-1}}
\]

le signe dépend de la convention d'orientation (qu'il faudra préciser le moment venu !).

D'autre part
\[
 \deleted{\operatorname{Tr}xyz=1}\qquad \tau-\tau'=\operatorname{Tr}\underbrace{[x,y]}_z.z=\operatorname{Tr}z^2=\operatorname{Tr}1=2
\]

d'où
\[
\begin{aligned}
\deleted{\unsure{1|2}}2&=2\varepsilon\varepsilon_0\qquad\text{i.e.}\\
\varepsilon&=\deleted{\tfrac12}\varepsilon_0\ \deleted{\unsure{\varepsilon}}
\end{aligned}
\]

d'où
\[
 \boxed{\tau-\tau'=\deleted{\tfrac12}\varepsilon_0\underbrace{(1\wedge x\wedge y\wedge z)\otimes e^{-1}}_{\Delta\ \mathrm{déf}}}
 \qquad\text{i.e.}\quad 2(\tau-\tau')=\varepsilon_0\Delta
\]
(a un sens si $2$ inv. dans l'anneau de base $k$ où on travaille).

On en tire
\[
 4\tau\tau'=(\tau+\tau')^2-(\tau-\tau')^2=\Sigma^2-\deleted{\tfrac14}\Delta^2
\]

À vérifier (dans $\mathbb Q[[\deleted{b,b',b''}\unsure{X_{ij}|x_{ij}},y_{ij},\unsure{Z_{ij}|z_{ij}}]]$ \quad $1\leq i,j\leq2$)
\[
 \Sigma^2-\deleted{\tfrac14}\Delta^2\in4\mathbb Z[b,b',b'',t,t',t'']
\]
\margin{?}
\deleted{i.e.}
\[
 \deleted{4\Sigma^2-\Delta^2\in16\mathbb Z[b,b',b'',t,t',t'']}
\]

d'où l'expression
\[
 \tau\tau'=\Pi(b,b',b'',t,t',t'')\quad\in\mathbb Z[b,b',b'',\unsure{t|\psi},t',t'']
\]
expression polynomiale de poids total $6$ (les $b$ étant de
% archive sheet number: 184
\scanpage{186}
On va plutôt essayer de calculer
\[
 \tau'-\tau=\operatorname{Tr}xyz-\operatorname{Tr}yxz=\operatorname{Tr}([x,y]z)=\operatorname{Tr}([y,z]x)=\operatorname{Tr}([z,x].y)
\]

C'est une expression \inserted{multilin.} alternée en $x,y,z$, donc on doit avoir
\[
 \tau'-\unsure{\tau|z}=\deleted{\unsure{\varepsilon|1/2|\operatorname{Tr}}}(a\wedge x\wedge y\wedge z)\otimes e^{-1}
\]
pour $a\in\unsure{A|\mathcal A}$ bien déterminé. Comme cette expression est invariante par $G=\underline{\mathrm{Aut}}\,A$, \deleted{on voit que} et que (si $g\in G(-)$)
\[
 g(a)\wedge g(x)\wedge g(y)\wedge g(z)=a\wedge x\wedge y\wedge z\qquad(=a\wedge g(x)\wedge g(y)\wedge g(z))
\]
on en conclut que $a=ga$, i.e. $a$ inv. sous $G$ i.e. $a=\varepsilon.1$. Donc on doit avoir
\[
 \tau'-\tau=\varepsilon(1\wedge x\wedge y\wedge z)\otimes e^{-1}
\]
et pour calculer $\varepsilon$, il suffit de prendre $x,y,z$ tels que $1\wedge x\wedge y\wedge z$ inv. i.e. $1,x,y,z$ une base \inserted{de $A$}. On prendra

\[
 x_0=\begin{pmatrix}0&0\\1&0\end{pmatrix},\ y_0=\begin{pmatrix}0&1\\0&0\end{pmatrix},\ z_0=[x_0,y_0]=\begin{pmatrix}-1&0\\0&1\end{pmatrix}=1-2\alpha
 \quad\text{où }\alpha=\begin{pmatrix}1&0\\0&0\end{pmatrix}
\]

(NB \(\displaystyle x_0y_0=\begin{pmatrix}0&0\\0&1\end{pmatrix},\ y_0x_0=\begin{pmatrix}1&0\\0&0\end{pmatrix}\))
\deleted{$x_0y_0z_0=\begin{pmatrix}0&0\\0&1\end{pmatrix}$)}

On trouve
\[
 1\wedge x_0\wedge y_0\wedge z_0=-2\,1\wedge x_0\wedge y_0\wedge\alpha
 =-2\alpha'\wedge x_0\wedge y_0\wedge\alpha
 =2\deleted{x\wedge y\wedge\alpha\wedge\alpha'}\inserted{\alpha\wedge\alpha'\wedge x\wedge y}
\]
où
\[
 \alpha=\begin{pmatrix}1&0\\0&0\end{pmatrix},\quad \alpha'=\begin{pmatrix}0&0\\0&1\end{pmatrix}
\]
et \deleted{$x,y$,} $\alpha,\alpha',x_0,y_0$ font base can. de $M_2(\mathbb Z)$. On
% archive sheet number: 185
\scanpage{187}
poids $1$ et les $t$ de poids $2$), et de poids total $2$ en chacune des \deleted{paires de} variables \deleted{$(b,t)\ (b',t')\ (b'',t'')$} \inserted{$x,y,z$ ($x$ intervenant avec le poids $1$ dans $b,t',t''$, poids $0$ dans $b',b'',t$)}. De plus, il doit être symétrique \deleted{\unsure{par}} permutation des \deleted{\unsure{quatre}} \inserted{trois} paires de variables $(b,t),(b',t'),(b'',t'')$. Et \unsure{itou|idem} pour $\Delta^2$.

On doit donc avoir
\[
\begin{aligned}
\Delta^2={}&\alpha tt't''+
\deleted{\beta(b^2t't''+t b^{\prime 2}t''+tt'b^{\prime\prime 2})+\gamma(t b^{\prime 2}b^{\prime\prime 2}+b^2t'b^{\prime\prime 2}+b^2b^{\prime 2}t'')+\delta b^2b^{\prime 2}b^{\prime\prime 2}}\\
&\deleted{\unsure{\Gamma|\gamma}(b'b''b^{\unsure{2}}t^{\unsure{\prime}}+b'b\,b^{\prime\prime 2}t^{\unsure{\prime}}}\\
&\inserted{\beta(b^2t^2+b^{\prime 2}t^{\prime 2}+b^{\prime\prime 2}t^{\prime\prime 2})+\gamma(b'b''t't''+b''bt''t+bb'tt')+}\\
&+\delta(b'b''b^2t+b''bb^{\prime 2}t'+bb'b^{\prime\prime 2}t'')+\unsure{\eta|\zeta}\, b^2b^{\prime 2}b^{\prime\prime 2}
\end{aligned}
\]
avec
\[
 \alpha,\beta,\gamma,\delta\in\mathbb Z
\]

à déterminer, par (quatre ?) choix particuliers judicieux des triples $(x,y,z)$. P. ex. le triple précédent $x_0,y_0,z_0=[x_0,y_0]$ donne $b,b',b''=0$, et $\Delta=2\varepsilon_0$, $\underline{\Delta^2=4}$, d'où
\[
 4=\alpha tt't''
\]
où
\[
 t''=\operatorname{Tr}x_0y_0=1,\quad t^{\deleted{\unsure{\prime}}}=\operatorname{Tr}y_0z_0=0,\quad t'=\operatorname{Tr}z_0x_0=0
\]

et on trouve $4=0.\alpha=0$, ce qui est fâcheux !\unsure{.|!} Et pour cause -- il faudra, dans l'expression de $\Delta^2$, faire intervenir aussi $c,c',c''$. En fait, on trouve immédiatement, une fois qu'on y pense :
% archive sheet number: 186
\scanpage{188}
\[
 \Delta(x,y,z)^2=\det\begin{pmatrix}
 \varphi(1,\deleted{\unsure{x}}1)&\varphi(1,x)&\varphi(1,y)&\varphi(1,z)\\
 \varphi(x,1)&\varphi(x,x)&\varphi(x,y)&\varphi(x,z)\\
 \varphi(y,1)&\varphi(y,x)&\varphi(y,y)&\varphi(y,z)\\
 \varphi(z,1)&\varphi(z,x)&\varphi(z,y)&\varphi(z,z)
 \end{pmatrix}
\]
où $\varphi(u,v)$ est la forme bil. sym. associée à la forme quadratique $q$ ($=\det$)
\[
 \varphi(u,v)=\operatorname{Tr}u\operatorname{Tr}v-\operatorname{Tr}(uv)=
\]

\uline{NB} On a
\[
 \begin{aligned}
 \varphi(u,1)&=\operatorname{Tr}(u)\deleted{\unsure{-\varphi(u,1z)|-\varphi(v,1z)}}\\
 \varphi(u,u)&=2\det u
 \end{aligned}
\]

Donc on trouve
\[
 \boxed{\Delta(x,y,z)^2=\det\begin{pmatrix}
 2&b&b'&b''\\
 b&2c&bb'-t''&bb''-t'\\
 b'&bb'-t''&2c'&b'b''-t\\
 b''&bb''-t'&b'b''-t&2c''
 \end{pmatrix}}
\]
ce qui permet d'expliciter
\[
 4\tau\tau'=\Sigma^2-\Delta^2
\]
\deleted{\unsure{où|$\longleftarrow$}}
où $\Delta$ est donné ci-dessus, et
\[
 \boxed{\Sigma=bt+b't'+b''t''-bb'b''}
\]
% archive sheet number: 187
\scanpage{189}
Le calcul ``bébête'', \unsure{en développant|du développ.t} le déterminant et tenant compte des symétries (pour regrouper les termes qui se correspondent par \deleted{symétrie} op. de $\unsure{\mathfrak S_3|\mathscr S_3}$) et $\Sigma^2$,
\[
 \Sigma^2-\Delta^2=4\Pi(b,b',b'',t,t',t'',c,c',c'')
\]
avec
\[
\boxed{\begin{aligned}
 \Pi={}&tt't''+(ct^2+c't^{\prime 2}+c''t^{\prime\prime 2})+(b^2c'c''+b^{\prime 2}c''c+b^{\prime\prime 2}cc')\\
 &-(b'b''ct+b''bc't'+bb'c''t'')
\end{aligned}}
\]
\[
 =tt't''+ct(t-b'b'')+c't'(t'-b''b)+c''t''(t''-bb')+(b^2c'c''+b^{\prime 2}c''c+b^{\prime\prime 2}cc')
\]

On aura donc
\[
 \left\{\begin{aligned}
 \tau+\tau'&=\Sigma&&\bigl(=bt+b't'+b''t''-bb'b''\bigr)\\
 \tau\tau'&=\Pi&&\bigl(=tt't''+(ct^{\unsure{2|\prime}}+c't^{\prime 2}+c''t^{\prime\prime 2})^{\unsure{,|\prime}}\\
 &&&\quad+(b^2c'c''+b^{\prime 2}c''c+b^{\prime\prime 2}cc')\\
 &&&\quad-(b'b''ct+b''bc't'+bb'c''t'')\bigr)
 \end{aligned}\right.
\]
\[
 \tau-\tau'=\varepsilon_0\Delta
\]
% archive sheet number: 188
\scanpage{190}
\uuline{À vérifier} : L'anneau \inserted{$B$} des fonctions $G$-invariantes \unsure{sur|\deleted{sur}} \inserted{$F(x,y,z)$ en trois variables dans} $\unsure{A|\mathcal A}$ (alg. d'Azumaya de rang $2$ sur $k$\unsure{)|\}}, $G=\underline{\mathrm{Aut}}\,(A)\simeq\underline{A}^{*}/\unsure{\mathbb G_m|G_m}$) est une algèbre \deleted{libre} \inserted{finie libre} de rang $2$ sur l'anneau des polynômes à neuf variables
\[
 A=k[\underline{\unsure{B|b}},\underline B',\underline B'',\underline{\unsure{C|c}},\underline C',\underline C'',\underline{t},\underline t',\underline t'']
\]
où
\[
 \left\{\begin{aligned}
 \underline B&=\operatorname{Tr}x,& \underline B'&=\operatorname{Tr}y,& \underline B''&=\operatorname{Tr}z\\
 \underline C&=\det x,& \underline C'&=\det y,& \underline C''&=\det y\\
 \underline t&=\operatorname{Tr}yz,& \underline t'&=\operatorname{Tr}zx,& \underline t''&=\operatorname{Tr}xy
 \end{aligned}\right.
\]

engendré par la fonction
\[
 \tau=\operatorname{Tr}xyz\qquad(=\operatorname{Tr}yzx=\operatorname{Tr}z^{\unsure{\prime}}xy)
\]
soumis à la relation quadratique
\[
 \boxed{\tau^2-\Sigma\tau+\Pi=0}
\]
où \deleted{$\Sigma$\unsure{ et}, $\Pi$}
\[
 \Sigma,\Pi\in A
\]

sont donnés plus haut. Si $u\mathbin{\unsure{-|\mapsto}}u'$ est l'involution canonique de
\[
 B\simeq A[\underline\tau]/(\underline\tau^2-\Sigma\underline\tau+\Pi),
\]
% archive sheet number: 189
\scanpage{191}
définie par
\[
 \tau\longrightarrow\tau'=\Sigma-\tau,
\]

on aura
\[
 \tau'=\operatorname{Tr}yxz=\operatorname{Tr}xzy=\operatorname{Tr}zyx.
\]

Les \unsure{élémts|éléments} du groupe symétrique $\unsure{\mathfrak S_3|\mathscr S_3}$ échangent entre eux les neuf générateurs $\underline b,\underline c,\underline t$, et aussi $\tau$ et $\tau'$ (les permutations circ. laissent $\tau,\tau'$ fixes, les transpositions échangent $\tau,\unsure{\tau''|\tau'}$).

En \deleted{termes} tant qu'anneau où $\mathfrak S_3$ opère, on peut décrire $B$ comme
% The curved transposition mark leads from the group of t's to the point after the double-primed c.
\[
 B=k[\underline b,\underline b',\underline b'',\underline c,\underline c',\underline c'',\inserted{\underline t,\underline t',\underline t'',}\underline\tau,\underline\tau']/
 \bigl(\underset{\substack{\text{équations séparé-}\\\text{ment invariantes}\\\text{sous }\mathfrak S_3}}{\underbrace{\underline\tau+\underline\tau'-\Sigma},\ \underbrace{\tau\tau'-\Pi}}\bigr)
\]

avec l'opération évidente de $\mathfrak S_3$ sur l'anneau des polynômes à onze variables $\underline b,\underline{\unsure{C|c}},\underline t,\underline\tau,\underline\tau'$.

\uline{NB} Quand $2$ est inv. dans $k$, et dans ce
% archive sheet number: 190
\scanpage{192}
cas seulement, on peut choisir comme générateur de $B$ sur $A$
\[
 \Delta=\tau-\tau'
\]
soumis à la relation
\[
 \Delta^2=\Sigma^2-4\Pi\quad\in A
\]

On aura
\[
 \begin{aligned}
 \tau&=\tfrac12(\Delta+\Sigma)\\
 \tau'&=\tfrac12(-\Delta+\Sigma)
 \end{aligned}
\]

Mais en car. 2, on aura
\[
 \Delta=\tau+\tau'=\Sigma\in A\qquad\text{(en car. 2)}
\]
donc $\Delta$ ne risque pas d'engendrer $B$ sur $A$ en car. 2, donc pas non plus si $2$ non inv. dans $k$ i.e. $k/2k\ne0$.

\begin{center}\rule{2.4cm}{0.3pt}\end{center}

Cas de quatre variables $\unsure{x_1,y_1,z_1,t_1|x,y,z,t}$.

Soit $B_4(\mathcal A)\subset\mathrm{Sym}(\overset{\vee}{\mathcal A}^{4})$ l'anneau des fonctions polynômes $F(x,y,z,t)$
% archive sheet number: 191
\scanpage{193}
invariantes par $G$. On y trouve les éléments
\[
 \left.\begin{aligned}
 b_i&=\operatorname{Tr}x_i\\
 c_i&=\det x_i
 \end{aligned}\right\}\qquad1\leq i\leq4
\]
\[
 t_{ij}=\deleted{\unsure{t}}t_{\{i,j\}}=\operatorname{Tr}x_ix_j\qquad\{i,j\}\in\unsure{\mathcal P_2|P_2|\mathfrak P_2}([1,4])
\]
(ce qui fait déjà $4+4+6=14$ fonctions sur un ``espace'' quotient'' de dim $4\times4-3=13$ -- il doit y avoir \deleted{\unsure{donc|deux}} \inserted{une} relation entre ces quantités), et enfin les
\[
 t_{ijk}=\operatorname{Tr}(x_ix_jx_k)\qquad
 \deleted{1\leq i\unsure{<|\leq}j\unsure{<|\leq}k\leq4}\quad
 \inserted{i,j,k\in[1,4]\ \text{tous différents}}
\]
qui sont algébriques sur l'anneau engendré par les $b_i,c_i,t_{ij}$.

Il y a $24$ fonctions $t_{ijk}$, \deleted{mais} (\unsure{correspondants|correspondant} aux \unsure{repères|copies} de l'ens à $4$ éléments) mais
\[
 t_{ijk}=t_{jki}=t_{kij}
\]
donc il n'y a que $\unsure{8|4}$ fonctions distinctes, lesquelles se groupent $2$ par $2$ \deleted{\unsure{en}}
\[
 \left\{\begin{aligned}
 t_{ijk}+t_{jik}&=\Sigma_{\{i,j,k\}}\in\deleted{\unsure{\mathbb Z|k}}\mathbb Z[b_i,c_i,t_i,b_j,c_j,t_j,b_k,c_k,t_k]\\
 t_{ijk}.t_{jik}&=\Pi_{\{i,j,k\}}\in\text{id}
 \end{aligned}\right.\tag*{$(*)$}
\]
% archive sheet number: 192
